Honest performance was robust to the analytic choices a developer could plausibly vary (Table~\ref{tab:e6}). On DEPRESJON, the full feature set gave AUROC 0.768 to 0.785 (logistic regression, XGBoost); dropping any single feature group changed it to 0.740 to 0.797; single groups alone ranged from 0.458 (day/night, logreg) to 0.802 (temporal/spectral, xgboost); SMOTE inside training folds gave 0.761 to 0.786 with ECE 0.097 to 0.109; leave-one-subject-out gave 0.745 to 0.773. On PSYKOSE, the full feature set gave AUROC 0.902 to 0.903 (logistic regression, XGBoost); dropping any single feature group changed it to 0.845 to 0.901; single groups alone ranged from 0.703 (day/night, logreg) to 0.875 (distributional, logreg); SMOTE inside training folds gave 0.904 to 0.906 with ECE 0.046 to 0.082; leave-one-subject-out gave 0.899 to 0.906. On HYPERAKTIV, the full feature set gave AUROC 0.452 to 0.466 (logistic regression, XGBoost); dropping any single feature group changed it to 0.418 to 0.493; single groups alone ranged from 0.459 (circadian, xgboost) to 0.519 (temporal/spectral, logreg); SMOTE inside training folds gave 0.454 to 0.466 with ECE 0.149 to 0.208; leave-one-subject-out gave 0.424 to 0.472. No variant on HYPERAKTIV rose above chance under subject-wise evaluation, and no variant on DEPRESJON or PSYKOSE approached the record-wise estimates of E1: the gap between published and honest numbers is a property of the splitting design, not of feature engineering, resampling or the cross-validation scheme.
